\section{GPT-3 y ChatGPT: cómo los impactos externos reescribieron el criterio de la empresa}

La inteligencia cognitiva marcó la dirección interna, pero la ola de los grandes modelos reescribió una y otra vez esa ruta desde fuera. GPT-3 y ChatGPT fueron dos impactos de distinto nivel: uno cambió la imaginación técnica; el otro, a los usuarios y al mercado.

Que la entrevista trate a GPT-3 y a ChatGPT en capítulos separados es importante. GPT-3 representa la entrada del paradigma de escala en el horizonte; ChatGPT representa el momento en que la percepción de los usuarios y el ritmo de la industria se encendieron de golpe. Cuando Zhang Peng (张鹏) recuerda estos momentos, lo hace con un tono «a la vez ansioso y emocionado»: la ansiedad venía de que la ola externa era demasiado grande; la emoción, de que la dirección para la que se habían preparado durante tanto tiempo por fin quedaba validada.

\begin{figure}[H]
\centering
\includegraphics[width=0.96\textwidth]{model-shockwaves.png}
\caption{Tres impactos de modelos: GPT-3, ChatGPT, DeepSeek y los cambios de paradigma posteriores.}
\end{figure}

\begin{knowledgebox}{Lectura de la figura: cada impacto cambia un nivel}
GPT-3 cambió la imaginación técnica: la escala podría producir capacidades generales. ChatGPT cambió la percepción de los usuarios: cualquier persona supo lo que un gran modelo puede hacer. DeepSeek cambió el discurso sobre eficiencia y código abierto: la capacidad de un modelo no depende solo de invertir recursos masivos; también puede impulsarse con ingeniería y con un ecosistema abierto.
\end{knowledgebox}

\subsection{Por qué «la oportunidad es para quien está preparado»}

Después de los impactos externos, esta sección responde a un malentendido frecuente: si aprovechar una oportunidad depende solo de la suerte. La respuesta de Zhang Peng se acerca más a la acumulación cotidiana que a una predicción prodigiosa.

Zhang Peng cita a su mentor: una oportunidad es como un tablón en el mar; cuando llega flotando, todavía tienes que dar unas cuantas brazadas para alcanzarlo. La ola de los grandes modelos es muy difícil de predecir con precisión, pero sí es posible prepararse a largo plazo: mantener la dirección correcta y acumular sin pausa capacidades de ingeniería, de investigación y de comercialización, para tener con qué aprovechar la oportunidad cuando llegue.

\begin{importantbox}{Prepararse no es predecir}
Predecir es saber qué ocurrirá mañana; prepararse es seguir acumulando capacidades en la dirección correcta aunque no se sepa qué ocurrirá mañana. La ola de los grandes modelos se parece más a lo segundo.
\end{importantbox}

\subsection{Resumen de la sección}

La ruta de modelos de Zhipu (智谱) no consistió en seguir sin más a GPT-3 o a ChatGPT, sino en que su dirección temprana de inteligencia cognitiva fue corregida repetidamente por los paradigmas técnicos externos. Cuanto mayor es el impacto externo, más se pone a prueba si la organización está preparada.
